\chapter{Fazit}\label{Fazit}
\section{Erfüllung der Anforderungen}
In diesem Kapitel folgt die Beurteilung, ob die in \chapref{Anforderungen Tool} gestellten Anforderungen unter Berücksichtigung der Ergebnisse von \chapref{chap:result_benchmark} durch das Verfahren erfüllt werden konnten. 

\begin{itemize}
    \item \textbf{Die Wahl der verwendeten Testframework(s) ist unerheblich}\\
    Das Verfahren ist frameworkunabhängig für jedes Testframework, welches die in \chapref{chap:Anforderungen_Framework} gestellten Anforderungen erfüllt, also für jedes Framework, welches JUnit XML als Ausgabeformat bereitstellt, und eine eindeutige Zuordnung zwischen Testausgabe und Testcase über die Attribute \texttt{name} und \texttt{classname} ermöglicht.
    \item \textbf{Nach Abschluss des Verfahrens wurde jeder Testcase für den aktuellen Stand des Codes mindestens einmal durchgeführt}\\
    Die Implementation von \texttt{testAuditor} in \chapref{Komponente_testAuditor} garantiert eine korrekte Ermittlung der fehlenden Testergebnisse, da nach \chapref{chap:limit_scope} die Korrektheit des Inputs angenommen wird. 
    \item \textbf{Die Ausgabe des durch das Verfahren ausgeführten Testlaufs ist auf dieselbe Weise verfügbar, wie die anderen Ausgaben}\\
    Diese Anforderung kann von einem korrekt implementierten \texttt{controller} erfüllt werden, da die \texttt{testExecution} auf dasselbe Testframework wie die vorherigen Testläufe zugreift, und ihr Ergebnis an den \texttt{controller} zurückmeldet. Die genau Vorgehensweise ist zwar von der jeweiligen Implementation des \texttt{archive} und \texttt{controllers} anhängig, die notwendigen Voraussetzungen zur Erfüllung dieser Anforderung sind jedoch durch das Verfahren erfüllt.
    \item \textbf{Die Laufzeit eines Testdurchlaufs unter Verwendung des Verfahrens ist kleiner oder gleich der Laufzeit eines konventionellen Durchlaufs}\\
    Diese Anforderung konnte mit Blick auf \chapref{chap:summary_benchmark} nur teilweise erfüllt werden, da die Laufzeit für einen Durchlauf ohne Inhalte im \texttt{archive} langsamer als ein konventioneller Durchlauf ist. Jedoch wird die Laufzeit des Verfahrens kürzer, je mehr Testcases bereits ausgeführt wurden und deren Ergebnis im \texttt{archive} liegt. Der Punkt, ab dem das Verfahren schneller abläuft, als ein konventioneller Durchlauf, ist abhängig von Anzahl und Ausführungsgeschwindigkeit der Testcases.
\end{itemize}

\section{Schlussbetrachtung}
Zusammenfassend lässt sich sagen, dass durch das in dieser Bachelorarbeit entwickelte Verfahren durchaus eine Verkürzung der Laufzeit von Testläufen erzielen lässt. Es wurde erklärt, von welchen Faktoren der genau Wert der Verkürzung in einem Projekt abhängig ist. 

Außerdem wurden die Schwierigkeiten beim Versuch, eine möglichst umfassende Unabhängigkeit vom benutzten Testframework zu erzielen, dargestellt. Es wurde die breite Verwendbarkeit des Verfahrens in verschiedenen Testframeworks, aber auch die Limitationen dieser Unabhängigkeit genau durchleuchtet. 

Die zu Beginn gestellte Frage:\\
\textit{\textbf{(Wie) Kann die Laufzeit von Testläufen unabhängig vom verwendeten Testframework durch die Wiederverwendung bereits vorhandener Testergebnisse reduziert werden, ohne die Aussagekraft der Testergebnisse zu beeinträchtigen?}}\\
kann also mit \textit{\textbf{''Ja, unter gewissen Voraussetzungen.''}} beantwortet werden. Die Klärung des \textit{\textbf{''Wie''}} ist Inhalt von \chapref{Vergleich Testausgabefomate} und \chapref{chap:approach}. 

\section{Möglichkeiten zur Verbesserung}
In seiner aktuellen Form kann das Verfahren nicht überprüfen, ob die bereitgestellten XML-Ausgaben der Tests auch tatsächlich aus dem aktuellen Stand des Codes entstammen, oder ob im Nachhinein Änderungen an dem zu testenden Code oder den Tests selber vorgenommen wurden. Die Verwendung veralteter Ergebnisse kann die Aussagekraft des Verfahrens erheblich beeinträchtigen.\\
Eine mögliche Lösung wäre, jedes Testergebnisses eindeutig zu einem Commit zuzuordnen. So könnte zumindest überprüft werden, ob das Ergebnis aus dem aktuellen Commit entstammt.

Ein anderes Problem ist, dass das Tool alle Tests, deren Ergebnis nicht im \texttt{archive} vorhanden ist, als nicht abgelaufen. Allerdings ist es möglich, dass Testcases in einer vorherigen Iteration ausgeführt wurden, und seither sich weder die Testcases noch der betroffene Code geändert hat. In diesem Fall kann das Ergbnis dieser Testcases weiterhin als valide betrachtet werden, und diese Testcases müssen nicht erneut ausgeführt werden.\\
Eine Möglichkeit, hier nachzubessern, wäre die Integration eines neuen Tools, welches das hier entwickelte Verfahren mit einer Regression Test Selection wie in \chapref{RTS} verbindet. So kann die Menge der Testcases, die abgelaufen werden müssen, weiter beschränkt werden, was eine zusätzliche Verringerung der Laufzeit verspricht.  

Eine weitere Möglichkeit, das Verfahren zu verbessern, besteht in der Implementation von zusätzlichen Parsern in \texttt{testAuditor} für andere Formate als JUnit XML, wie beispielsweise TAP oder OTR. Dies würde die Integration der Verfahrens in Projekte ermöglichen, die eines diese weiteren Formate, aber kein JUnit XML nutzen. 

Zusätzlich könnte untersucht werden, ob eine effizientere Methode der Kommunikation mit dem Testframework als mit dem aktuellen Aufbau möglich ist, um das Verfahren bei bereits geringerer Abdeckung des \texttt{archive} effizienter als einen konventionellen Durchlauf zu machen. 
